\documentclass[12pt]{article}
\usepackage[utf8]{inputenc}
\usepackage[T1]{fontenc}
\usepackage{hyperref}
\title{Tarao Naga Language Resources for K-12 Curriculum Development}
\author{Compiled by Claude Sonnet 5 (claude-sonnet-5) for Andr\'as Kornai}
\date{}
\begin{document}
\maketitle

\section{Introduction}
Tarao (ISO 639-3 tro) is a small language of Chandel and Ukhrul districts, Manipur, India, with about 870 people counted in the 2001 census. Wikipedia classes it as Kuki-Chin, and UNESCO lists it as critically endangered. Very little material was found.

\section{Primary Educational Resources}
No primers, textbooks or school programmes were found.

\section{Linguistic and Reference Materials}
\begin{itemize}
\item Language overview (Wikipedia), which cites a Tarao grammar by Chunghkham Yashwanta Singh (2002, New Delhi: Akansha): \url{https://en.wikipedia.org/wiki/Tarao_language}
\item Tarao people (Wikipedia): \url{https://en.wikipedia.org/wiki/Tarao_people}
\item Endangered Languages Project profile and bibliography: \url{https://endangeredlanguages.com/lang/3504} and \url{https://endangeredlanguages.com/lang/3504/bibliography}
\end{itemize}
Caution: a ``Dictionary of the Taraon Language'' (1963) appears in catalogues, but Glottolog identifies it as Digaru (Digaro-Mishmi), a different language. It is not listed as a Tarao resource.

\section{Religious and Community Resources}
No Tarao Bible translation or community organisation was found in the search results.

\section{Digital and Audio Resources}
None were found.

\section{Curriculum Development Context}
The reference description says Tarao is not known to be taught in schools, and nothing found contradicts this. A curriculum would need primary documentation work with the community first.
\end{document}